\documentclass[11pt]{article}
\usepackage[margin=1in]{geometry}
\usepackage{fontspec}
\setmainfont{TeX Gyre Termes}
\newfontfamily\cjkfont[BoldFont={AR PL Mingti2L Big5}]{AR PL Mingti2L Big5}
\usepackage{amsmath,amssymb,mathtools}
\usepackage{booktabs}
\usepackage{enumitem}
\usepackage{microtype}
\usepackage{xcolor}
\usepackage[numbers,sort&compress]{natbib}
\usepackage[hidelinks]{hyperref}
\providecommand{\tightlist}{\setlength{\itemsep}{0pt}\setlength{\parskip}{0pt}}

\newenvironment{zhabstract}
  {\begin{quotation}\cjkfont\small\raggedright\noindent\textbf{繁體中文摘要}\par\smallskip\mdseries}
  {\end{quotation}}

\title{Canonicalization Before Quadratic Audit:\\
A Certificate-Methods Architecture for an Oriented Level-11 Owner Ledger}
\author{Liang Wang\textsuperscript{1}\\
\small \textsuperscript{1}School of Artificial Intelligence and Automation,\\
\small Huazhong University of Science and Technology,\\
\small Luoyu Road 1037, 430070, Hubei, P.R. China\\
\small \texttt{wangliang.f@gmail.com}}
\date{2 September 2026}

\begin{document}
\maketitle

\begin{abstract}

This article revises the research architecture for a finite ownership
problem associated with a positive time change of the
\texttt{Gamma\_0(11)} geodesic flow. The frozen population contains 138
Hecke-output instances in 55 source-word/prime groups, with oriented
primitive \texttt{Gamma\_0(11)} conjugacy classes as the proposed
owners. The earlier design made 9,453 terminal pair dispositions appear
to be the foundational certificate. Adversarial review identified that
requirement as an unjustified architectural lock. Revision 1 instead
makes a deterministic canonicalization map and its biconditional the
primary target: two rooted, oriented inputs must receive the same
canonical owner bytes exactly when they represent the same oriented
primitive owner. The 9,453-row table is retained, but only as a derived
adversarial audit of the canonical partition. Aggregate class counts
remain weaker post-closure controls. The article also defines the
distinct roles of a global owner table \texttt{G}, the 138-row incidence
relation \texttt{I}, and the cell-local no-double-credit quotient
\texttt{C}. The executed work is limited to closed-corpus evidence
synthesis and review-informed method design. No canonicalization
theorem, owner partition, all-pairs audit, or \texttt{G/I/C} table has
been produced. All 22 inherited citations retain \texttt{anchor:none};
source identities and metadata close, but claim-to-passage faithfulness
remains \texttt{INCONCLUSIVE}. The contribution is therefore a
reproducible, fail-closed certificate architecture rather than a
scientific ownership result. P31 remains A1-only, with no formal Route-A
tuple, positive arithmetic A2 credit, or Route-B invocation.
\end{abstract}

\begin{zhabstract}
本文針對固定的 $\Gamma_0(11)$ 測地流正時間變換， 提出一套失敗即關閉的 所有權憑證方法架構。 凍結輸入包含 一百三十八個 Hecke 輸出實例、 五十五個來源字與質數群組； 預定所有者是 保持方向的本原 $\Gamma_0(11)$ 共軛類。 核心目標不是先製作 九千四百五十三筆 彼此獨立的成對判定， 而是建立可重播的 決定性規範化映射及其雙條件： 兩個已取本原根 且保留方向的輸入 具有相同規範位元， 當且僅當它們代表 同一所有者。 只有在此定理、 證據格式與獨立驗證器完成後， 全部成對表格 才可作為派生的 對抗式回歸稽核。 本文並嚴格區分 全域所有者表 $G$、 一百三十八列的關聯表 $I$ 與儲存格內去重商 $C$。 本階段僅執行 封閉語料的文獻綜合 與審查後方法設計， 未實作規範化器、 共軛判定、 本原根證書、 所有者分割 或任何 $G/I/C$ 結果。 所有引文均無段落定位， 逐主張的來源支持仍屬未定； 本文不宣稱新定理、 普查結果、 新穎性或路線晉級。
\end{zhabstract}

\noindent\textbf{Keywords:} canonical forms; modular geodesic flow;
$\Gamma_0(11)$; oriented conjugacy; primitive owners; certificate methods;
adversarial audit

\medskip
\noindent{\cjkfont\textbf{中文關鍵詞：}規範形式；模群測地流；保持方向共軛；本原所有者；憑證方法；對抗式稽核}

\hypertarget{1-introduction-research-question-and-contribution}{%
\section{Introduction, Research Question, and
Contribution}\label{1-introduction-research-question-and-contribution}}

Finite orbit ledgers require an equivalence rule before they can support
a dynamical interpretation. A repeated row may represent the same
primitive owner, a traversal power, an inverse-oriented owner, or a
recurrence of one owner in another correspondence cell. P31 freezes
those possibilities rather than allowing an outcome-dependent merge. Its
inherited dynamical object is the positive time-changed flow
\texttt{X\_geo/rho\_epsilon} on \texttt{T\^{}1Y\_0(11)}. The real
level-11 newform differential, positivity interval, Hecke normalization,
reciprocal log-zeta convention, and period coordinate \texttt{k=2y+z}
remain unchanged.

The finite input is also unchanged: 138 instances are distributed over
55 source-word/prime groups. The inherited \texttt{2/2/134} split and
the three 55-group diagnostic summaries are instance- or group-level
controls, not owner counts. The proposed owner is an oriented primitive
conjugacy class in \texttt{Gamma\_0(11)} represented by a positive-trace
determinant-one lift. Inversion is kept as a separate oriented owner and
connected by an inverse-link field. A primitive traversal exponent is
not a Hecke branch-cycle degree. Equality of trace, length, homology, or
any other filter cannot by itself certify subgroup conjugacy.

The revised research question is:

\textbf{Can a deterministic, independently replayable canonicalization
contract certify the complete oriented primitive-owner partition of the
frozen 138-instance population, and, conditional on that closure, induce
the distinct global, incidence, and cell-local estimands \texttt{G},
\texttt{I}, and \texttt{C}?}

This formulation changes the method priority without changing the object
or scientific state. The primary target is a canonicalization
biconditional, not 9,453 independently foundational negative
certificates. The all-pairs expansion remains valuable because it can
expose inconsistencies, nontransitive implementation defects,
inverse-policy mistakes, and compensating merges or splits. It is
nevertheless derived from the partition certificate unless a later
theorem proves that bespoke pairwise evidence is necessary for a
particular disposition.

The article's contribution is a certificate-methods architecture with
three separated surfaces: a canonical proof object, a full-population
adversarial audit, and downstream estimands. It is not positioned as a
new conjugacy theorem, an implemented solver, or a novelty claim. No
novelty search was authorized in the frozen research program. The article type is therefore
deliberately narrower than a mathematical-results paper.

Three distinctions govern the exposition. First, the dynamical owner is
not an input row: ownership is assigned only after subgroup conjugacy,
orientation, and maximal-root status are resolved. Second, a proof
object is not an estimand: canonical bytes and their witnesses justify a
partition, whereas \texttt{G}, \texttt{I}, and \texttt{C} summarize
different consequences of that partition. Third, an audit is not its
own truth source: the all-pairs expansion can challenge a canonicalizer,
but it cannot define correctness when the canonicalizer's biconditional
is still open. Keeping these levels separate makes a negative or
not-evaluable outcome reportable without silently changing the owner
definition.

This architecture also fixes the temporal order of evidence. The owner
contract, theorem versions, serialization, and failure states must be
registered before the 138 rows are classified. The pair table may be
generated only after those inputs have independently replayed
certificates. Downstream counts may be read only after the pair audit
agrees with the canonical partition or records a typed discrepancy.
Nothing in this order presumes that the required canonicalizer exists;
the lawful endpoint may remain \texttt{NOT\_EVALUABLE}.

\hypertarget{2-frozen-literature-and-theoretical-boundary}{%
\section{Frozen Literature and Theoretical
Boundary}\label{2-frozen-literature-and-theoretical-boundary}}

\hypertarget{21-exact-subgroup-representations}{%
\subsection{Exact subgroup
representations}\label{21-exact-subgroup-representations}}

The frozen literature describes modular-subgroup structure through
several exact mathematical interfaces.
Subgroup classification, arithmetic-geometric descriptions, special
polygons, and two congruence-recognition routes supply foundational
finite-index and finite-data contexts
\citep{P31-S01,P31-S02,P31-S03,P31-S04,P31-S05}.
% ARS-CITE source_ids=P31-S01,P31-S02,P31-S03,P31-S04,P31-S05 anchor=none claim_to_passage=INCONCLUSIVE
A later computational fundamental-domain framework supplies a potential
exact Fuchsian-domain interface
\citep{P31-S06}.
% ARS-CITE source_ids=P31-S06 anchor=none claim_to_passage=INCONCLUSIVE

These sources support only a bounded synthesis finding: finite,
replayable subgroup data are plausible ingredients. They do not bind the
exact P31 representation, prove the required canonicalization
biconditional, decide any frozen pair, or define owner bytes.
Representation, decision, and certification remain distinct obligations.

For the present article, this is a negative-boundary finding rather than a claim of
literature absence. The corpus was frozen by a bounded search and has no
passage-level locators. We can say that no retained source was admitted
as the complete project solver; we cannot say that no such theorem or
method exists anywhere. A later implementation would need a source-
finalization step that records the exact representation hypotheses and
the passages licensing each subroutine. Until then, even a mathematically
promising route is an unbound component rather than an available
certificate engine.

\hypertarget{22-ambient-and-arithmetic-conjugacy-components}{%
\subsection{Ambient and arithmetic conjugacy
components}\label{22-ambient-and-arithmetic-conjugacy-components}}

The second source group gives candidate components for matrix and
arithmetic conjugacy. Ideal-class/matrix correspondences and their later
refinements supply ambient integral structure, including a relation
between hyperbolic integral matrices and ideal classes
\citep{P31-S07,P31-S08,P31-S09}.
% ARS-CITE source_ids=P31-S07,P31-S08,P31-S09 anchor=none claim_to_passage=INCONCLUSIVE
Continued-fraction, arithmetic-group, and modern integral-matrix methods
enlarge the possible ambient decision toolkit
\citep{P31-S10,P31-S11,P31-S12,P31-S13}.
% ARS-CITE source_ids=P31-S10,P31-S11,P31-S12,P31-S13 anchor=none claim_to_passage=INCONCLUSIVE

No source in this group may be promoted into the required oriented
\texttt{Gamma\_0(11)} solver. A valid specialization would still have to
preserve determinant one, positive-trace lifting, the congruence
condition, orientation, primitive roots, inverse linkage, termination,
and independently replayable evidence. Failure of a bounded ambient
search would not constitute a negative subgroup-conjugacy certificate.

Ambient-to-subgroup transfer is a particularly important stop surface.
Two matrices may be related under an ambient group while the required
conjugator fails the $\Gamma_0(11)$ condition, or an ambient procedure
may omit the orientation and root policies needed by the ledger. The
future dossier must therefore identify the decision group at every
step, prove that normalizations remain within it, and serialize the
conjugating or obstructing evidence. Merely reporting that a general
algorithm terminated would not show that the project-specific owner
predicate was evaluated.

\hypertarget{23-canonical-reduction-roots-inversion-and-replay}{%
\subsection{Canonical reduction, roots, inversion, and
replay}\label{23-canonical-reduction-roots-inversion-and-replay}}

Canonicalization requires more than an ambient yes/no decision.
Modular-geodesic reduction provides relevant coding context
\citep{P31-S15}.
% ARS-CITE source_ids=P31-S15 anchor=none claim_to_passage=INCONCLUSIVE
Word-hyperbolic conjugacy algorithms supply individual and batch decision
precedents
\citep{P31-S16,P31-S17}.
% ARS-CITE source_ids=P31-S16,P31-S17 anchor=none claim_to_passage=INCONCLUSIVE
Centralizers and reversing symmetries clarify why ordinary symmetry and
inversion must be typed separately
\citep{P31-S18}.
% ARS-CITE source_ids=P31-S18 anchor=none claim_to_passage=INCONCLUSIVE
Pell equations, unit methods, and computational algebraic number theory
provide potential arithmetic subroutines and implementation background
\citep{P31-S19,P31-S20}.
% ARS-CITE source_ids=P31-S19,P31-S20 anchor=none claim_to_passage=INCONCLUSIVE

The synthesis does not choose among polygon, arithmetic, or
hyperbolic-group implementations. It instead specifies the acceptance
condition that any route must meet. A route is admissible only if its
hypotheses are proved for the frozen marked object, its output is
canonical under the orientation policy, its primitive-root and inverse
fields are exact, and an independent verifier can replay success and
failure evidence. The literature does not supply the project
serialization or prove that this combined contract terminates.

\hypertarget{24-aggregate-counts-as-post-closure-controls}{%
\subsection{Aggregate counts as post-closure
controls}\label{24-aggregate-counts-as-post-closure-controls}}

Class-number structure supplies useful aggregate context
\citep{P31-S14}.
% ARS-CITE source_ids=P31-S14 anchor=none claim_to_passage=INCONCLUSIVE
A direct formula for primitive hyperbolic class counts in
\texttt{Gamma\_0(N)} is the closest frozen census precedent, and modular
conjugacy has also been related to real-quadratic class numbers
\citep{P31-S21,P31-S22}.
% ARS-CITE source_ids=P31-S21,P31-S22 anchor=none claim_to_passage=INCONCLUSIVE

The information content of an aggregate count is lower than that of a
labeled partition. Compensating false merges and false splits can
preserve a total. Class counts can therefore test a closed owner table
for consistency, but they cannot choose canonical representatives,
certify an individual owner identity, or establish that every input has
been resolved.

\hypertarget{3-executed-methodology}{%
\section{Executed Methodology}\label{3-executed-methodology}}

\hypertarget{31-closed-corpus-literature-synthesis}{%
\subsection{Closed-corpus literature
synthesis}\label{31-closed-corpus-literature-synthesis}}

The only executed research method in this phase is literature synthesis
over frozen artifacts. The upstream corpus process captured 44 records,
removed nine duplicate manifestations, screened 35 unique records,
excluded 13, and retained 22 sources. The inclusion frame covered
modular-subgroup descriptions, integral and subgroup conjugacy,
reduction and centralizers, primitive roots, and direct class-count
precedents. Nineteen inventory rows are classified as peer-reviewed. The
authorized metadata correction fixed the P31-S16 page range to 287--305.

The source-verification layer closed identity and metadata for 22 of 22
rows within its recorded scope. That layer mainly used DOI, publisher,
journal, institutional, metadata, abstract, and limited
authoritative-record surfaces. It did not inspect every theorem passage,
produce claim-level page or section locators, or perform a general
current retraction or source-conflict screen. The evidence matrix then
assigned each source an admitted contribution, an excluded stronger
claim, and an applicability warning. Phase 3 grouped the sources into
representation, conjugacy, canonicalization, and census themes without
converting proximity into theorem transfer.

The synthesis used source-effect discipline rather than vote counting.
A source could support a subgroup representation or an ambient decision
method while being explicitly excluded from proving the composite P31
owner contract. Publication status and venue recognition did not erase
that applicability boundary. Similarly, several adjacent algorithms do
not become a complete solver merely because their advertised tasks can
be listed in the same paragraph. The missing theorem-to-certificate
links remain visible as method obligations.

Citation closure was checked structurally: every source identifier used
by the manuscript belongs to the frozen inventory and resolves to one
bibliography entry. That check does not validate a theorem passage. All
citations retain \texttt{anchor:none}, no direct quotation is used, and
claim-to-passage status remains \texttt{INCONCLUSIVE}. The wording is
therefore deliberately limited to component context, prospective use,
and recorded exclusions.

\hypertarget{32-review-adjudicated-revision-procedure}{%
\subsection{Review-adjudicated revision
procedure}\label{32-review-adjudicated-revision-procedure}}

Phase 5 supplied four procedurally separated, single-model-family review
records: editorial, ethics, citation-integrity, and Devil's Advocate.
Their integrated decision was \texttt{MAJOR\_REVISION}, with no Critical
finding and no ethics \texttt{BLOCKED} result. Revision 1 applies the
author-adjudicated branch recorded in the Phase-6 contract. It changes
study design only: canonicalization becomes primary, the all-pairs
ledger becomes a derived audit, \texttt{G/I/C} receive self-contained
conditional definitions, and the AI disclosure is narrowed to the actual
verification surface.

No new retrieval occurred during Revision 1. No source field, reference,
passage locator, direct quotation, experiment, code result, proof
result, or canonical manuscript byte was introduced. The Phase-6
ClaimIntent manifest was frozen before this prose and supplies the
complete set of eight substantive article claims. Every finding below is
either a closed-corpus evidence-synthesis statement, a project
definition, or a prospective method obligation.

\hypertarget{4-review-adjudicated-certificate-architecture}{%
\section{Review-Adjudicated Certificate
Architecture}\label{4-review-adjudicated-certificate-architecture}}

\hypertarget{41-primary-target-a-canonicalization-biconditional}{%
\subsection{Primary target: a canonicalization
biconditional}\label{41-primary-target-a-canonicalization-biconditional}}

Let \texttt{X} denote the frozen set of 138 instances after exact input
validation. The future method should define a partial operation
\texttt{root(x)} that either returns an oriented primitive
\texttt{Gamma\_0(11)} representative together with a traversal exponent
and proof payload, or returns a typed not-evaluable state. It should
then define a deterministic byte map

\begin{verbatim}
kappa: X -> OwnerBytes
\end{verbatim}

on every successfully resolved input. The primary theorem target is the
biconditional

\begin{verbatim}
kappa(x)=kappa(y)
if and only if
root(x) and root(y) represent the same oriented primitive Gamma_0(11) owner.
\end{verbatim}

This is a target, not a proved statement. Its forward direction must
prevent accidental byte collisions. Its reverse direction must show that
all representatives of one oriented owner reduce to identical bytes.
Both directions must be established in the exact subgroup and
presentation used by P31. The map must not identify an owner with its
inverse; instead, an \texttt{inverse\_owner\_bytes} field should connect
the two oriented objects when that relation is defined.

The biconditional identifies the mathematical certificate invariant: a
total, sound, complete, deterministic owner map. If such a theorem and
its per-instance witnesses are independently replayable, the partition
follows from byte equality. A separate bespoke negative proof for every
unequal pair is then not logically foundational merely because the
population is finite.

The two implications have different failure modes. Soundness fails if
two distinct oriented owners can collide in one byte string. Completeness
fails if two representatives of one owner can survive with different
bytes. Determinism fails if the same input can produce different bytes
under permitted executions. Totality fails if any validated instance
lacks either a certificate or a typed unresolved disposition. A future
theorem statement should name these obligations individually; a single
headline such as ``canonical reduction works'' would be too coarse for
the ledger contract.

The primitive-root layer must precede canonical owner serialization.
If an input is a traversal power, serializing the unreduced input could
mint a false owner. Conversely, root extraction must not erase the
traversal exponent needed by the incidence relation. The proposed
certificate therefore carries both the primitive owner bytes and the
exact repetition field. Hecke degree remains an input coordinate and is
never substituted for that repetition field.

\hypertarget{42-prospective-certificate-and-verifier-contract}{%
\subsection{Prospective certificate and verifier
contract}\label{42-prospective-certificate-and-verifier-contract}}

Each future per-instance certificate should bind at least: the immutable
input identifier and hash; the subgroup representation and theorem
version; exact membership evidence; the normalized matrix or word; the
maximal primitive root and traversal exponent; the orientation
convention; canonical owner bytes; inverse linkage; and the complete
proof or reduction trace needed by a read-only verifier. A failure
certificate should identify the exact failed precondition or unresolved
theorem obligation rather than silently convert a timeout into
nonconjugacy.

The verifier acceptance predicates are also prospective. It should
recheck input binding, determinant and subgroup membership, root
powering, primitiveness under the chosen theorem, orientation
preservation, deterministic serialization, and any inverse relation. It
should reject unknown fields, stale theorem versions, noncanonical
encodings, missing proof data, or an unresolved subroutine. Independent
replay means an implementation separate from the producer can evaluate
the frozen mathematical predicates; it does not mean that two AI reviews
have statistically independent errors.

No such schema, theorem binding, fixture set, producer, or verifier was
implemented in the recorded research or writing stages. The specification above partially addresses
reproducibility at the article level while leaving scientific and
implementation closure explicitly open.

The producer and verifier should communicate through a versioned,
closed schema. Required fields should have fixed encodings, and unknown
fields should cause rejection rather than permissive parsing. Each proof
payload should state which theorem and representation version it uses,
so that a change in a normal-form convention cannot silently reuse old
bytes. A manifest should bind the accepted input hashes, fixture hashes,
producer build, verifier build, and output hashes. These requirements
describe a prospective reproducibility contract; no such manifest or
build exists in the present article.

Positive and negative evidence also require asymmetric handling. A
positive conjugacy certificate can carry an explicit conjugator whose
membership and action are replayed. A negative disposition needs a
theorem-backed exhaustive argument appropriate to the frozen group and
representation; failure to find a conjugator is insufficient. Root
certificates face the same distinction: displaying a power relation is
positive evidence, while claiming no proper root requires an exhaustive
licensed method. The verifier must preserve these differences instead
of reducing every result to one Boolean field.

Target-blind fixtures should exercise the contract before population
execution. At minimum they should include identical inputs, distinct
representatives of one owner, inverse-related owners, proper powers,
coarse-invariant collisions, malformed membership data, and deliberately
unresolved theorem bindings. Expected outputs must be registered before
the producer sees them. A fixture suite can reveal implementation
defects, but passing it cannot replace the global biconditional or the
complete 138-instance replay.

\hypertarget{43-the-9453-row-table-as-a-derived-adversarial-audit}{%
\subsection{The 9,453-row table as a derived adversarial
audit}\label{43-the-9453-row-table-as-a-derived-adversarial-audit}}

For \texttt{\textbar{}X\textbar{}=138}, the unordered-pair expansion has
\texttt{binom(138,2)=9,453} rows. Once \texttt{kappa} is total, each row
can carry the two input IDs, their owner bytes, the derived equality
disposition, inverse relationship, and any optional direct-solver
cross-check. The expected same-owner relation is equality of canonical
bytes; inequality supplies the derived different-owner relation under
the proved biconditional.

This table remains useful. It can test symmetry, transitivity
consequences, deterministic sorting, collision handling, inverse-policy
consistency, and agreement between canonical and direct routes.
Adversarial fixtures can target pairs that share trace, length,
homology, or other coarse invariants while differing in exact ownership.
Aggregate-count controls can then operate on the completed partition
without replacing any row-level audit.

The table is not described as a uniquely necessary proof architecture.
If a future scientific requirement needs a separately replayable
negative obstruction for particular cross-class pairs, that need must be
justified and bound. Until then, the full table is a regression
expansion of the canonical certificate rather than 9,453 independent
foundations.

The audit should retain discrepancies rather than overwrite them. If a
direct pair solver and canonical-byte equality disagree, the row should
record both results, their proof payloads, and a fail-closed batch state.
Neither route should be selected after inspecting which one preserves a
preferred aggregate count. This makes the table an adversarial consumer
of the canonical partition rather than a mechanism for repairing the
partition post hoc.

Pair rows can also test equivalence-relation consequences globally. A
same-owner relation derived from canonical bytes is automatically
reflexive, symmetric, and transitive at the byte level, but producer or
serialization defects may still appear as inconsistent input bindings,
inverse labels, or traversal metadata. The 9,453-row expansion exposes
all unordered cross-input comparisons and therefore remains the most
demanding finite regression surface even though it is not the primary
proof object.

\hypertarget{44-distinct-g-i-and-c-estimands}{%
\subsection{\texorpdfstring{Distinct \texttt{G}, \texttt{I}, and
\texttt{C}
estimands}{Distinct G, I, and C estimands}}\label{44-distinct-g-i-and-c-estimands}}

The downstream objects can be defined conditionally once \texttt{kappa}
is total and the biconditional is proved. Define the global owner table

\begin{verbatim}
G = {kappa(x): x in X},
\end{verbatim}

with exactly one row per distinct oriented owner byte string. Define the
incidence relation \texttt{I} with one row for every input instance,
carrying its input fields, cell coordinates, owner bytes, traversal
exponent, inverse link, and all frozen Hecke coordinates. Thus
\texttt{\textbar{}I\textbar{}=138} even when
\texttt{\textbar{}G\textbar{}\textless{}138}.

Let a cell key be the frozen tuple
\texttt{(source\_word,\ prime,\ hecke\_degree)}. Define \texttt{C} as
the set of distinct \texttt{(cell\_key,\ owner\_bytes)} pairs induced by
\texttt{I}. Raw multiplicity remains visible in \texttt{I}, while
\texttt{C} awards one unit to an owner within a cell. The same owner may
appear once in each of several cells without being duplicated in
\texttt{G}.

Conditional well-definedness is straightforward but important. If
\texttt{kappa} is total and satisfies the biconditional, changing an
input representative within the same oriented owner cannot change its
\texttt{G} row or its cell-level owner key. Deduplication inside a cell
is therefore invariant under representative choice. Conversely, two
different canonical owner bytes cannot be merged in \texttt{C} without
violating the definition. This set-theoretic argument does not prove
that \texttt{kappa} exists or that the frozen inputs are resolved; it
states what follows if the primary certificate closes.

The three estimands answer noninterchangeable questions. The cardinality
of \texttt{G} concerns distinct global oriented owners represented in
the population. The rows of \texttt{I} retain provenance and allow one
owner to be traced back to every input occurrence. The rows of
\texttt{C} enforce no-double-credit within a frozen correspondence cell
while allowing the same owner to occur in several cells. Publishing only
one of these surfaces would destroy information needed to reconstruct
the other two. For that reason their schemas, validation rules, and
summary statistics should be declared independently.

Materialization must remain conditional on zero unresolved owner rows.
If even one input cannot be rooted or canonicalized, \texttt{I} may
retain a diagnostic record but \texttt{G} and \texttt{C} cannot be
reported as complete estimands. A partial table may be useful for method
development only if it is labeled incomplete and does not award owner
counts. This stop rule prevents missing inputs from being silently
treated as new, different, or irrelevant owners.

\hypertarget{5-evidence-synthesis-findings}{%
\section{Evidence-Synthesis
Findings}\label{5-evidence-synthesis-findings}}

Eight precommitted positions control this manuscript. First, the frozen
corpus supplies complementary subgroup, ambient-conjugacy, arithmetic,
reduction, and hyperbolic-group components, but no executed complete
oriented \texttt{Gamma\_0(11)} owner solver. Second, the primary
certificate target is the deterministic canonicalization biconditional.
Third, the 9,453-row table is a derived adversarial audit, and aggregate
counts remain post-closure controls. Fourth, \texttt{G}, \texttt{I}, and
\texttt{C} are distinct estimands populated only after owner closure.

Fifth, a complete prospective certificate preserves subgroup membership,
orientation, primitive-root status, inverse linkage, traversal exponent,
deterministic serialization, and replayable failure evidence while
keeping Hecke degree separate. Sixth, the article's methodological
contribution is the separation of canonical proof object, audit
expansion, and estimands. Seventh, metadata and source-identity closure
do not clear theorem passages, retraction status, source conflicts, or
the theorem-to-certificate bridge. Eighth, the manuscript remains A1-only
and changes no scientific result or Route state.

These are evidence-synthesis and design findings. They do not imply
feasibility, novelty, theorem correctness, or implementation readiness.
The central scientific state remains \texttt{NOT\_EXECUTED}.

\hypertarget{6-reproducibility-and-prospective-interface}{%
\section{Reproducibility and Prospective
Interface}\label{6-reproducibility-and-prospective-interface}}

Reproducibility in the present article means that another reader can
recover the frozen corpus, Phase-4 report, Phase-5 reviews, Phase-6
manifest, and revision log from named hash-bound artifacts. The
literature selection and source ledger are inspectable, and the
citation/reference/source-ID sets close. Generative prose is not
promised to be byte-reproducible, and source-passage judgments cannot be
replayed because the locators were never finalized.

A later scientific package should freeze four layers before reading
owner outcomes. Layer one binds the exact \texttt{Gamma\_0(11)}
representation and proves every input conversion. Layer two binds the
canonicalization and root theorem, including both directions of the
biconditional. Layer three freezes producer and verifier schemas,
adversarial fixtures, and typed failure states. Layer four defines the
derived all-pairs audit and the conditional \texttt{G/I/C}
materialization. Hashes, theorem versions, serialization versions, and
fixture manifests should be immutable across execution.

The smallest valid next test is not the census. It is a
theorem-to-certificate dossier evaluated against deliberately small,
target-blind fixtures. Scientific population execution requires separate
authorization after the primary biconditional, root policy, inverse
policy, and independent verifier have closed. A missing theorem or
unresolved fixture should terminate as
\texttt{NOT\_EVALUABLE\_CONJUGACY\_INCOMPLETE}, not as a negative owner
result.

\hypertarget{7-discussion-and-implications}{%
\section{Discussion and
Implications}\label{7-discussion-and-implications}}

The revised architecture improves the match between the mathematical
object and its evidence. A partition is naturally represented by a total
equivalence invariant. Expanding that invariant to every pair can be an
excellent audit, but expansion size should not be confused with proof
content. This distinction reduces a false quadratic necessity while
preserving the most demanding adversarial check available for the fixed
population.

The change does not make P31 easy. The canonicalization theorem must
simultaneously handle exact subgroup membership, oriented conjugacy,
primitive roots, inverse linkage, and deterministic bytes. A compact
certificate that lacks either direction of the biconditional would be
weaker than the all-pairs design it replaces. The revision therefore
narrows the foundational object without lowering the exactness standard.

The \texttt{G/I/C} separation also clarifies what a future result would
mean. \texttt{G} answers how many global oriented primitive owners
occur. \texttt{I} records how the 138 frozen instances map to those
owners. \texttt{C} answers how owner identity contributes within each
frozen correspondence cell. None of the inherited instance summaries can
substitute for these estimands, and no cell-local recurrence can
multiply the number of global owners.

The architecture is intentionally compatible with more than one future
mathematical implementation. A polygonal reducer, an arithmetic route,
or a word-hyperbolic route could be considered, provided its exact
hypotheses and evidence are bound to the same owner semantics. This
interoperability is conditional: different producers cannot be compared
until they preserve orientation, root status, inverse linkage, and
serialization under one independently checked contract. The article
therefore standardizes the meaning of a prospective result without
claiming that any implementation route is presently adequate.

For Route A, this remains primitive-owner infrastructure at A1. No
primitive Euler product, dynamical determinant, zero comparison,
analytic continuation, or arithmetic specificity test is produced.
Editorial improvement cannot promote a Route coordinate. Route B remains
closed because there is no Route-A readiness, Hilbert space, operator,
domain, self-adjointness result, trace formula, or divisor identity.

\hypertarget{8-acknowledged-limitations}{%
\section{Acknowledged
Limitations}\label{8-acknowledged-limitations}}

The dominant limitation is citation resolution. All 22 prose citation
pairs retain \texttt{anchor:none}. The source ledger supports identity,
metadata, and bounded claim-fitness accounting, but the exact passages
and theorem hypotheses were not frozen. Claim-to-passage faithfulness is
therefore \texttt{INCONCLUSIVE}. Revision 1 narrows language rather than
inventing locators.

The corpus has no general current retraction or source-conflict
clearance. Those checks remain unrun and their status remains unknown.
The P31-S16 page correction is
preserved, but one corrected field does not constitute a clean integrity
screen. The corpus is historically weighted and no new contribution or
novelty comparison was authorized.

The method architecture is uninstantiated. There is no bound
canonicalization theorem, no producer or independent verifier, no
completed fixture suite, no owner decision, and no \texttt{G/I/C}
output. The conditional definitions show how outputs would relate after
closure; they do not close the primary mathematical obligation.

Finally, AI-assisted review and drafting used one Codex model family
with procedural role separation. This supplies multiple documented
perspectives, not statistically independent validation. Liang Wang's
gate confirmations establish author decisions and workflow authority
only; they are not evidence of personal full-text checking.

\hypertarget{9-future-work}{%
\section{Future Work}\label{9-future-work}}

Future work should proceed in this order:

\begin{enumerate}
\def\labelenumi{\arabic{enumi}.}
\tightlist
\item
  perform a separately authorized source-finalization pass that freezes
  exact theorem passages, hypotheses, correction status, and any
  required source-conflict or retraction checks;
\item
  bind one exact subgroup representation and prove the conversion from
  every frozen input;
\item
  state and prove the canonicalization biconditional, maximal-root
  contract, orientation convention, and inverse-link rule;
\item
  freeze certificate schemas, producer/verifier roles, typed failures,
  and target-blind adversarial fixtures;
\item
  execute the 138 per-instance certifications only after the interface
  recheck passes;
\item
  derive the 9,453-row audit and compare it with optional direct pair
  checks without treating aggregate counts as pair evidence; and
\item
  materialize and independently validate \texttt{G}, \texttt{I}, and
  \texttt{C} only after zero unresolved owner inputs.
\end{enumerate}

Each step requires a new gate when it extends beyond literature-only
manuscript composition. Nothing in this article authorizes retrieval or
scientific execution.

\hypertarget{10-conclusion}{%
\section{Conclusion}\label{10-conclusion}}

P31 now has a more defensible certificate architecture. The foundational
target is a total deterministic canonical owner map with a proved
biconditional for oriented primitive \texttt{Gamma\_0(11)} ownership.
The 9,453-row table remains a full-population adversarial audit derived
from that certificate, not a presumed uniquely necessary collection of
bespoke proofs. Aggregate class counts remain weaker controls, and the
global owner table \texttt{G}, incidence relation \texttt{I}, and
cell-local quotient \texttt{C} remain distinct conditional outputs.

This is concrete design-level progress, but it is not a scientific
result. No canonical form, owner partition, pair audit, theorem,
computation, or estimand table has been executed. The source corpus
remains passage-unresolved, originality remains unassessed, the formal
Route-A tuple remains \texttt{UNASSIGNED}, positive arithmetic A2
remains absent, and Route B remains closed.

\section*{Declarations}

\paragraph{Author contribution and accountability.}
Liang Wang specified the mathematical object and restrictions, supplied
the conceptual direction, approved the recorded stage gates,
adjudicated the certificate-first design, and is responsible for the
manuscript's scholarly content. AI assistance is disclosed below and is
not credited with authorship. Stage-gate approval does not attest that
the author personally performed full-text or exact source-passage
verification.

\paragraph{Funding.} No funding was received for this work.

\paragraph{Competing interests.} The author declares no competing
interests.

\paragraph{Ethics statement.} Human-subjects and animal-research review
is not applicable. The work concerns theoretical mathematics,
published-source records, and project-owned workflow artifacts; it
involves no participants, identifiable personal data, animals,
recruitment, or intervention.

\paragraph{Data and materials availability.} The materials are the
frozen 22-row source inventory and the hash-bound project artifacts
recorded for Round 10. No owner ledger, scientific dataset, solver
output, pair-decision table, or \texttt{G/I/C} result was generated during
manuscript preparation.

\hypertarget{ai-disclosure-and-verification-limitation}{%
\section*{AI Disclosure and Verification
Limitation}\label{ai-disclosure-and-verification-limitation}}

OpenAI Codex, using the GPT-5 model family, assisted during the session
dated 2026-09-02 UTC; the exact backend snapshot/build was not exposed.
AI-assisted work in the recorded pipeline included literature-search
support, source-identity and metadata checking, evidence-matrix
construction, evidence synthesis, report drafting, four role-based
Phase-5 reviews, review synthesis, ClaimIntent-constrained Revision-1
drafting, citation-ID/reference closure checks, and revision-log
accounting. No AI system executed a P31 solver, proof, pair decision,
owner census, experiment, or canonical-results refresh.

Liang Wang is the responsible human author. He approved the project
restrictions, stage gates, and the Phase-6 author-adjudicated design
choice. Those approvals must not be interpreted as a statement that he
personally read every source in full or verified any claim at the exact
source-passage level. The recorded verification was bounded mainly to
source identity, metadata, abstracts, authoritative landing pages, and
project-local claim-fitness records. All 22 citations lack passage
locators, so claim-to-passage faithfulness remains
\texttt{INCONCLUSIVE}; the article does not claim theorem-level source
verification, novelty clearance, or a clean retraction/conflict screen.

\bibliographystyle{plainnat}
\bibliography{references}

\end{document}
